\linksection{capturing}{Capturing fixtures}
\begin{description}
    \item[{\mono{capsys}}~] Capturing stdout/stderr in a test by overriding \py{sys.stdout},
          e.g.\ for \py[escapeinside=||]{print(|\codeskip*|)}.
    \item[\mono{capfd}~] \textsb{Capturing stdout/stderr in a test at file descriptor level,}
          e.g.\ for subprocesses or libraries printing from C code. \\[.5em]
          Type for both: \monot{pytest.CaptureFixture[\textcolor{pygkw}{str}]}
    \item[{\parbox[t]{6.5em}{\mono{capsysbinary}/\\\mono{capfdbinary}}}]
          \parbox[t]{6cm}{
              Like \mono{capsys}/\mono{capfd}, \\ but for binary data. \\
              Type: \monot{pytest.CaptureFixture[\textcolor{pygkw}{bytes}]}
          }
    \item[\mono{caplog}~] \textsb{Capturing logging output from a test.} \\
          Type: \mono{pytest.LogCaptureFixture}
\end{description}
\vspace{-.5em}

\mono{capsys} vs. \mono{capfd}: If in doubt, use \mono{capfd} to ensure everything is captured
(built-in pytest capturing also uses \mono{fd} by default).

\begin{minted}[autogobble, escapeinside=||]{python}
    def test_capfd(capfd):
        subprocess.run(["ls"])
        out, err = capfd.readouterr()
        assert out == "..."
\end{minted}

\vspace{-.5em}
Every process has \emph{two} outputs: \mono{stdout} and \mono[termred]{stderr}:
\vspace{-.5em}

\begin{minted}[autogobble,escapeinside=||]{text}
    $ ls doesnotexist /home
    |\textcolor{termred}{ls: cannot access 'doesnotexist': \codeskip}|
    |\textcolor{codeblue}{/home:}|
    |\textcolor{codeblue}{username}|
\end{minted}

\begin{minted}[autogobble,escapeinside=||]{text}
    $ ls doesnotexist /home |\highlightbox{> /dev/null}|
    |\textcolor{termred}{ls: cannot access 'doesnotexist': \codeskip}|
\end{minted}
\vspace{-.5em}

We get them separately with \mono{capfd.readouterr()}:

\begin{tabular}{l>{\columncolor{highlightcolor}}l>{\columncolor{highlightcolorred}}l}
                                       & stdout       & \highlightbox[highlightcolorred]{stderr} \\[.75em]
    \py{print("out 1")}                & \mono{out 1} &                                          \\
    \py{print("out 2")}                & \mono{out 2} &                                          \\
    \py{print("err", file=sys.stderr)} &              & \mono{err}                               \\
\end{tabular}

\vspace{-.25em}
\hspace{.2em} \py{out, err = capfd.readouterr()}  \\[.25em]
\hspace{1em} $\Rightarrow$ \hspace{.5em} \highlightbox{\mono{out}}: \py{"out 1\nout2\n"} \\[-.13em]
\hspace{1em} \phantom{$\Rightarrow$} \hspace{.5em} \highlightbox[highlightcolorred]{\monot{err\vphantom{t}}}: \py{"err\n"} \\[1em]

\mono{caplog} lets you access logging records. The \mono{.records} attribute provides full records which can be inconvenient (timestamps, etc.),
convenience attributes are available:

\vspace{-.5em}
\begin{minted}[autogobble,highlightlines={3,7-10}]{python}
def test_msgs(caplog: pytest.LogCaptureFixture):
    logging.warning("some msg")
    assert caplog.messages == ["some msg"]

def test_tpl(caplog: pytest.LogCaptureFixture):
    logging.warning("some msg")
    assert caplog.record_tuples == [
        # logger, loglevel,       message
        ("root", logging.WARNING, "some msg"),
    ]
\end{minted}

\pagebreak
